\documentclass{article}

\usepackage{neurips_2026}

\usepackage[utf8]{inputenc}
\usepackage[T1]{fontenc}
\usepackage{hyperref}
\usepackage{url}
\usepackage{booktabs}
\usepackage{amsfonts}
\usepackage{amsmath}
\usepackage{amssymb}
\usepackage{nicefrac}
\usepackage{microtype}
\usepackage{xcolor}
\usepackage{graphicx}
\usepackage{subcaption}
\usepackage{float}

\usepackage{tabularx}
\usepackage{multirow}
\usepackage{adjustbox}

\usepackage{tikz}

\newcommand{\R}{\mathbb{R}}
\newcommand{\tzh}[1]{{\color{ForestGreen}{[Zhenghan: #1]}}}
\DeclareMathOperator*{\argmin}{argmin}

\title{Why Language Models Transfer to Time Series}

\author{%
  Anonymous Author(s)
}

\begin{document}

\maketitle

\begin{abstract}

\end{abstract}

% 0. Intro
% 1. Related work
%      0+1: 2p
% 2. Training setup
%      Experiments and main results
%      Transfer works at scale - benchmark results + comparisons to SotA
%      1.5p
% 3. Pretrained models affinity to transfer (2p)
%      LLM Vs LLM+TS (where is transfer helping)
% 4. What changes during finetuning (2p)
%      LLM+TS Vs TS (pretrain Vs random init)
% 5. Final how transfer works? (1.5p)
%      Conclusion \& Discussion \& Limitations

i. intro \& RW
ii. establish transfer actually happens
iii. backbone exists in language (existing manifold)
iv. finetuning projects existing manifold
    1. lora vs io vs full: only a low rank of the weights need to change (adapting the existing manifold)
    2. effective rank of each model shows FT collapses the PT manifold into the rich directions useful for TS
    3. FT uses a lower dimensional manifold than RI which is more complicated (sine wave circles)
    4. explaining what the manifold represents using the crosscoder
v. Conclusion \& Discussion \& Limitations



% ═══════════════════════════════════════════════════════════════
% The phenomenon: LLMs finetuned on time series match or beat domain-specific models.
% The question: Why should representations learned from text help with temporal data?
% We provide a concrete, mechanistic account.
%
% Our approach: investigate from first principles.
%   1. Establish that transfer actually happens and quantify it. -> raw benchmark perf of pretrain Vs random
%   2. Characterize what the pretrained model already has (what does it contain that is useful for time series?) Analysis on QWEN BASE as it pertains to transfer 
%   3. Characterize what changes during finetuning (what changes in the representation space?) Analsysi on QWEN finetuned Compare with the control (learning from scratch) comparison of QWEN finetuend w random
%   4. Conclusion on how transfer works 
%
% Preview the main finding: pretraining creates a geometric prior —
% a low-dimensional, smooth representation space with directions
% that naturally encode temporal dynamics. Transfer = selecting those directions.

\section{Introduction}
\label{sec:intro}

% The key contributions we make in this paper are as follows:
% \begin{itemize}[leftmargin=*]
%   \item TODO
% \end{itemize} 
% ═══════════════════════════════════════════════════════════════

% ═══════════════════════════════════════════════════════════════
% LLM-to-TS transfer literature
% See Alexis...bib, section TRANSFER for papers
% Share manifold hyptohesis
%   shared embedding paper? 
% Any works on cross domain representational analysis 

\section{Related work}
\label{sec:related_work}

\paragraph{Time-series foundation models (TSFMs).}
Models designed as general-purpose time series forecasters, trained on Large-scale multi-domain data
\begin{enumerate}
    \item Chronos \citep{}
    \item Lag-Llama \citep{}
    \item MOMENT \citep{}
    \item TimeGPT \citep{}
\end{enumerate}

\paragraph{LLM-to-time-series transfer.} 
Work that repurposes pretrained language models for time series, primary related work
\begin{enumerate}
    \item Are Language Models Actually Useful for Time Series Forecasting?" \citep{tan2024languagemodelsactuallyuseful}
    \item Time-LLM: Time Series Forecasting by Reprogramming Large Language Models \citep{jin2024timellmtimeseriesforecasting}
    \item Using Pre-trained LLMs for Multivariate Time Series Forecasting \citep{wolff2025usingpretrainedllmsmultivariate}
    \item isionTS: Visual Masked Autoencoders Are Free-Lunch Zero-Shot Time Series Forecasters \citep{chen2024visionts}
\end{enumerate}

\paragraph{Time Series Tokenization and Representation; LLM2TS continued}
\begin{enumerate}
    \item A Time Series is Worth 64 Words \citep{nie2023timeseriesworth64}
    \item Small Vocabularies, Big Gains \citep{roger2025smallvocabulariesbiggains}
\end{enumerate}

\paragraph{The Platonic Representation Hypothesis? Cross-Modal Geometry?}
Theoretical/empirical work on representation convergence to explain why might representations transfer across modalities -> Language to TS
\begin{enumerate}
    \item The Platonic Representation Hypothesis \citep{huh2024platonicrepresentationhypothesis}
    \item Harnessing the Universal Geometry of Embeddings \citep{jha2025harnessinguniversalgeometryembeddings}
    \item Training Large Language Models to Reason in a Continuous Latent Space \citep{hao2024traininglargelanguagemodels}
\end{enumerate}

\paragraph{Scaling Laws and Efficient Adaptation}
How (do we) do models scale and fine-tuning strategies interact with transfer
\begin{enumerate}
    \item Scaling Laws for Transfer \citep{hernandez2021scalinglawstransfer}
    \item Studies on when scaling hurts? \citep{}
    \item Parameter-efficient fine-tuning comparisons, LoRA vs Full vs Adapters \citep{}
\end{enumerate}
The findings that pretrained 0.6B matches random 4B implies a transfer multiplier of ~8x in effective parameters; io\_only regime (embedding + lm\_head only) indicate if the transfer is fully captured by the i/o layers → PEFT/LoRA
% ═══════════════════════════════════════════════════════════════

% ═══════════════════════════════════════════════════════════════
\section{Establishing Transfer Learning is Real}
\label{sec:transfer}

% Setup paragraph: models + data + tokenization
% TODO: Prateek/Alexis to add benchmark details

\paragraph{Setup.} All models share the Qwen3 architecture. We train pairs of \textbf{FT} (pretrained $\to$ finetuned on TS) and \textbf{RI} (random init $\to$ trained on TS) across multiple scales (0.6B, 1.7B, 4B) using the same TS training procedure (uniform binning, 512 bins, context length 512). All models are evaluated on GiftEval (46 datasets).

\paragraph{Main result.} FT consistently outperforms RI across scales and datasets.

% ── TABLE: Benchmark results across scales ──
% TODO: Fill with actual ND[0.5] numbers from Alexis's sweeps
\begin{table}[h]
\centering
\caption{ND[0.5] on GiftEval (lower is better). FT outperforms RI at all scales.}
\label{tab:benchmark}
\begin{tabular}{lcccc}
\toprule
Model & 0.6B & 1.7B & 4B \\
\midrule
FT (pretrained + TS)  & TODO & TODO & TODO \\
RI (random + TS)      & TODO & TODO & TODO \\
\midrule
$\Delta$ (FT $-$ RI)  & TODO & TODO & TODO \\
\bottomrule
\end{tabular}
\end{table}

\paragraph{Scaling interaction.} Pretrained 0.6B matches or exceeds random-init 4B, implying a transfer multiplier of $\sim$8$\times$ in effective parameters.

% ── FIGURE: Performance across scales ──
% TODO: Bar chart or line plot showing FT vs RI at each scale
% Source: Alexis's training runs

\paragraph{Finetuning regime comparison.} We compare full finetuning, LoRA, and IO-only (embedding + LM head only).
% TODO: Table showing that even IO-only captures substantial transfer
% This previews the low-rank story in Section~\ref{sec:finetuning}

% ═══════════════════════════════════════════════════════════════

% ═══════════════════════════════════════════════════════════════
\section{Pretrained models affinity to transfer}
\label{sec:prior}
% ═══════════════════════════════════════════════════════════════

% Before any finetuning, does PT already have TS-compatible representations?
%
% 4.1 Linear decoding experiment
%   - A single linear projection from WikiText hidden states → TS
%   - No paired supervision, no model finetuning
%   - Result: 686 unique matches out of 10,000 real TS, 14/20 clusters
%   - This is only possible if the representation space already contains TS-like structure
% 4.1.1 Data distribution alignment with TS 
%   - Residual stream: h_{t+1} = h_t + f(h_t) → smooth trajectories
%   - Pretraining compresses to low-rank subspace
%   - Low-rank smooth trajectory projected to 1D = smooth 1D signal that resembles time series data distribution
%   - The distribution of projected trajectories overlaps with real TS distributions
%   - math to show that a model with representations which predict text sequences can easily finetune to predict time series data
% 4.1.2 The 2×2 ablation (architecture × input)
%   - text+PT: 686 unique — needs BOTH
%   - rand+PT: 4 unique — manifold exists but not traversed
%   - text+RandInit: 54 unique — variation but no manifold
%   - rand+RandInit: 99 unique
%   - Disentangles architecture from weights from input
%
% 4.2 Characterizing the prior: low dimensional geometry (transition to next section on geometry alignment)
%   - PCA: text+PT has effective rank ~6 in 28,672 dimensions
%   - RandInit: effective rank 431 (isotropic, no structure)

% ═══════════════════════════════════════════════════════════════
\section{What changes during finetuning}
\label{sec:finetuning}
% ═══════════════════════════════════════════════════════════════

% What does finetuning actually change in the representation space?
%
% 5.1 Subspace alignment
%   - can talk about LORA and linear finetuning results here
%   - ts@PT vs ts@FT: 60.6° (random baseline 85.3°) — most aligned pair while ts@PT vs ts@RI: 73.7° — much less aligned
%   - FT reuses PT's directions. RI finds different ones.
%
% 5.2 Weight changes (alexis will add)
%
% 5.3 Dimensional compression
%   - PT processes TS with eff rank 44
%   - FT compresses to eff rank 12 (found the useful directions)
%   - RI finds a different solution. eff rank 35 (higher dimensional, different subspace)
%   - FT achieves better compression BECAUSE it had the prior to select from
%   - Catestrophic forgetting: Text→FT collapses to rank 2.5 (from 150) 
%
% 5.4 Clean geometry vs complex reuse
%   - On periodic inputs (sine waves):
%     - RI: clean circular trajectories — the minimal representation, learned from scratch
%     - FT: topologically equivalent loops but geometrically more complex
%     - FT uses FEWER PCA dimensions because periodicity is projected onto
%       directions that were already high-variance from pretraining
%     - FT's loops carry the complexity of their linguistic origin
%   - This is direct evidence of reuse: FT repurposes, RI invents
%   - Ideally we show a convergence proof to the circle under random init which would help prove that transfer is really happening
%
% 5.5 Convergent features (FT_RI)
%   - Despite different geometry, FT and RI discover some of the same features
%   - Spike detectors, trough detectors, oscillation detectors appear in both
%   - These are the features that temporal data REQUIRES regardless of initialization
%   - But FT activates them at higher magnitude (the prior amplifies)
%   - Importantly how such features are implemented is different in FT and RI

% ═══════════════════════════════════════════════════════════════
\section{What Transfers?}
\label{sec:features}
% ═══════════════════════════════════════════════════════════════

% Zoom in: which specific neural features carry the transfer?
%
% 7.1 Crosscoder method
%   - Shared encoder applied independently per domain, per-domain decoders
%   - TopK=64 sparsity, 4096 latent dimensions
%   - Trained on 130k windows per layer, 27 layers
%   - WikiText probing with wiki-specific normalization
%
% 7.2 Feature categorization across layers
%   - PT_FT features (language transfer) concentrated in early-to-mid layers
%   - FT_RI features (convergent) dominant everywhere
%   - Category distribution changes with depth
%
% 7.3 Cross-domain features (PT_FT)
%   - Repetition detector (L8): "3·3·3·3" in math ↔ square waves in TS
%   - Slope detector (L26): "derivative", "tangent" ↔ rising/falling slopes
%   - Metronome detector (L4): "tempo", "beats" ↔ periodic spikes
%   - Chord progression detector (L5): harmonic sequences ↔ periodic peaks
%
% 7.4 The transfer is structural, not semantic
%   - Features detect computational patterns: repetition, gradients, periodicity
%   - These are primitives useful for ANY sequential prediction task
%   - Language pretraining happens to learn them because text also requires
%     tracking repetition, change, and rhythm

% ═══════════════════════════════════════════════════════════════
\section{Discussion and Limitations}
\label{sec:discussion}

\paragraph{The full picture.}
% Autoregressive pretraining creates smooth, low-dimensional trajectories (Section 3).
% This manifold contains directions naturally compatible with temporal data.
% Finetuning identifies the useful directions and compresses to them (Section 4.2),
% reusing pretrained structure rather than learning from scratch (Section 4.3).
% The transferred features are structural primitives—repetition, slopes, periodicity (Section 4.4).

\paragraph{Implications.}
% - Transfer should work for any pretrained autoregressive model, not just LLMs
% - The key property is smooth, low-dimensional sequential representations
% - The scaling result (pretrained 0.6B ≈ random 4B) suggests transfer provides
%   an ~8× effective parameter multiplier
% - LoRA's success suggests principled ways to preserve more pretrained structure

\paragraph{Limitations.}
% - Single architecture family (Qwen3), though multiple scales
% - GiftEval benchmark (diverse but finite)
% - Crosscoder dead features (40-60% per layer), limiting feature-level conclusions
% - The geometric analysis is descriptive—we characterize what changes but
%   cannot fully explain why specific directions are useful
% - Circuit analysis was inconclusive on the exact text computation being repurposed


\section{Conclusion}
\label{sec:conclusion}

% Language-to-TS transfer works because pretraining creates a geometric prior:
% a low-dimensional, smooth, anisotropic representation space.
% Finetuning projects this manifold onto time series by selecting the subset
% of pretrained directions useful for temporal data.
% The transferred features are structural primitives shared between language
% and time series—repetition, slopes, periodicity.
% The connection is geometric and structural, not semantic.

% ═══════════════════════════════════════════════════════════════

\end{document}
